%%\section{The Goods Market: Aggregate Demand and the Multiplier}

%%% [moved from MA_5_General_Equlibrium.tex, frame 1]
\begin{frame}{Say's Law}
    \mytext{
    Every demand equals supply, as we produce to have something to exchange for what we desire.

    Every supply equals demand, as everything produced will at some point be exchanged (or self-consumed).

    Production is the source of demand.
    }
\end{frame}

%%% [moved from MA_1_Fundamentals.tex, frame 23]
 

\begin{frame}{Keynesian Cross}
    \showmytikz{../shared/tikz/tikz-keynesiancross}
    \mycomment{
    Aggregate demand $AD$ is consumption, investment, and government expense;
    \begin{equation}
        AD = C+I+G
    \end{equation}
    with
    \begin{equation}
        C=a+b\left(Y-t\right)
    \end{equation}
    and $I,G,t$ constant; so that
    \begin{equation}
        AD = e + b\cdot Y
    \end{equation}
    with $G$ part of the autonomous part $e$.
    }
\end{frame}

%%% [moved from MA_1_Fundamentals.tex, frame 24]
\begin{frame}{Keynesian Cross}
    \showmytikz[2]{../shared/tikz/tikz-keynesiancross}
    \mycomment{
    An increase of government spending increases aggregate demand, shifting the $AD$ curve up. The result is a higher output $Y$.

    The same applies for autonomous increase of $I$.
    }
\end{frame}

%%% [moved from MA_1_Fundamentals.tex, frame 25]
\begin{frame}{Fiscal Multiplier}
    \mycomment{
    The idea with the multiplier is that any extra inserted expenditure (by fiscal policy, but anything else for that matter) is someone else's income, who will then re-spend a part of that, saving another part. 

    The size of the multiplier will depend on the propensity to spend (the inverse of the propensity to save), and on taxes paid each time.
    }
    \showmytikz{../shared/tikz/tikz-multiplier}
\end{frame}

%%% [moved from MA_1_Fundamentals.tex, frame 26]
\begin{frame}{Fiscal Multiplier}
    \mycomment{
    For now, let us abstract from taxes and net exports.
    }
    \mytext{
    Consumption $C$ is a (constant) part of income $Y$:
    \begin{equation}
        C=C_0+c\cdot Y
    \end{equation}
    Thus
    \begin{equation}
        Y=C_0+c\cdot Y + I + G
    \end{equation}
    For (in the short run) constant $I,C_0$ we get
    \begin{equation}
        \Delta Y = \frac{\Delta G}{1-c}=\frac{1}{s}
    \end{equation}
    with $s$ the savings quota. --
    Thus the multplier is
    \begin{equation}
        M=\frac{1}{1-c}=\frac{1}{s}.
    \end{equation}
    }
    \mycomment{
    % There are a number of (simplifying) assumptions here. One is that net exports $NX=Ex-Im$ be zero. Sin
    }
\end{frame}

%%% [moved from MA_1_Fundamentals.tex, frame 27]
\begin{frame}{Fiscal Multiplier: Critique}
    \mytext{
    \begin{itemize}
        \item crowding out private demand
        \item government debt means higher taxes:
        \item[$\rightarrow$] subjects reduce spending
        \item debt financing raises interest rates
        \item monetary policy counteracts
        \item lock-in problem
        \item might depend on government debt level
    \end{itemize}
    }
\end{frame}


%%% TODO (new content): automatic stabilisers; tax multiplier / balanced-budget multiplier; open-economy multiplier (import leakage)

